\documentclass[10pt,letterpaper]{article}

% Packages
\usepackage[utf8]{inputenc}
\usepackage{amsmath,amssymb}
\usepackage{booktabs}      % nice tables
\usepackage{graphicx}      % figures
\usepackage{hyperref}      % clickable refs
\usepackage[margin=1in]{geometry}
\usepackage{natbib}         % citations

% Title block
\title{[Paper Title]}
\author{[Authors]}
\date{}

\begin{document}
\maketitle

\begin{abstract}
% [CC fills this last]
\end{abstract}

\section{Introduction}
% [Background, motivation, contributions]

\section{Related Work}
% [Prior work organized by theme]

\section{Method}
% [Proposed approach with equations and algorithm]

\section{Experiments}

\subsection{Experimental Setup}
% [Dataset, metrics, baselines, implementation details]

\subsection{Main Results}
% [Table comparing baseline vs proposed]

% Example table:
% \begin{table}[h]
% \centering
% \caption{Main results on [Dataset].}
% \begin{tabular}{lcc}
% \toprule
% Method & Accuracy & F1 \\
% \midrule
% Baseline & 0.85 & 0.82 \\
% Proposed & \textbf{0.89} & \textbf{0.87} \\
% \bottomrule
% \end{tabular}
% \end{table}

% Example figure:
% \begin{figure}[h]
% \centering
% \includegraphics[width=0.8\textwidth]{figures/example.pdf}
% \caption{Overview of the proposed method.}
% \label{fig:overview}
% \end{figure}

\subsection{Ablation Study}
% [Remove each component to verify contribution]

\section{Conclusion}
% [Summary, limitations, future work]

\bibliographystyle{plainnat}
\bibliography{references}

\end{document}
